\bookchapter{Bab 7: Analisis lengkap fitur terkini}

\chapterintro{Satu, Komunikasi Satelit: Meminta Bantuan Tanpa Sinyal}

\booksection{Komunikasi satelit adalah terobosan besar di bidang ponsel dalam beberapa tahun terakhir. Saat Anda berada di daerah terpencil, hutan belantara, atau pegunungan dalam, di mana ponsel sama sekali tidak memiliki sinyal, komunikasi satelit bisa menjadi penyelamat.}

Prinsip Kerja:

\booksubsection{Cara Kerja:}

Model yang Didukung:

\booksubsection{Apple}

\begin{itemize}
  \item \textbf{Apple}: Seri iPhone 14 dan setelahnya, mendukung SOS darurat via satelit. Seri iPhone 17 lebih lanjut mendukung pengiriman iMessage dan SMS via satelit, dan memperluas cakupan ke lebih banyak wilayah.
  \item \textbf{Huawei}Xiaomi/OPPO/vivo
  \item \textbf{: Sejak 2024 model unggulan secara bertahap mendukung fitur pesan satelit BeiDou.}Tren perkembangan
  \item \textbf{: Komunikasi satelit sedang berkembang dari fitur eksklusif flagship kelas atas ke kelas bawah, di masa depan model kelas menengah juga diharapkan dapat mendukung. Sistem satelit BeiDou Tiongkok memiliki dukungan terluas di ponsel domestik.}Skenario Penggunaan:
\end{itemize}

\booksubsection{Petualangan outdoor, pendakian yang mengalami kecelakaan.}

\begin{itemize}
  \item Kecelakaan di laut.
  \item Bencana alam yang menyebabkan komunikasi darat terputus.
  \item Cara Penggunaan:
\end{itemize}

\booksubsection{Pastikan berada di area terbuka yang dapat melihat langit.}

\begin{enumerate}
  \item Buka fitur komunikasi satelit ponsel (biasanya di antarmuka panggilan darurat).
  \item Ikuti petunjuk layar, arahkan ke arah satelit.
  \item Kirim pesan bantuan dan lokasi yang sudah ditentukan.
  \item Pengetahuan Singkat
\end{enumerate}

\begin{tcbinfo}{Komunikasi satelit memerlukan ponsel dengan antena perangkat keras khusus, tidak semua ponsel mendukungnya. Selain itu, komunikasi satelit biasanya memerlukan layanan berlangganan berbayar, namun fitur panggilan darurat umumnya gratis.}
Dua, Komunikasi Tanpa Jaringan: Cara Baru Komunikasi Jarak Dekat

\end{tcbinfo}

\booksection{“Komunikasi tanpa jaringan” mengacu pada teknologi di mana dua ponsel dapat berkomunikasi langsung tanpa jaringan seluler dan Wi-Fi.}

Prinsip Teknologi:

\booksubsection{Prinsip Teknis:}

Skenario Aplikasi:

\booksubsection{Headphone nirkabel}

\begin{itemize}
  \item \textbf{: Komunikasi antar anggota tim saat mendaki atau berkemah.}Penyelamatan bencana
  \item \textbf{: Saat komunikasi terputus setelah gempa bumi atau banjir.}Acara besar
  \item \textbf{: Saat konser atau acara olahraga dengan banyak orang dan sinyal buruk.}Platform yang Didukung:
\end{itemize}

\booksubsection{Apple}

\begin{itemize}
  \item \textbf{Apple}Huawei
  \item \textbf{Huawei}Xiaomi
  \item \textbf{Xiaomi}Tiga, Deteksi Kecelakaan Mobil: Penyelamat di Saat Kritis
\end{itemize}

\booksection{Deteksi kecelakaan mobil adalah fitur yang menggunakan sensor ponsel untuk mendeteksi kecelakaan parah dan secara otomatis memanggil telepon darurat.}

Prinsip Kerja:

\booksubsection{Cara Kerja:}

Menampilkan antarmuka hitungan mundur, menanyakan apakah Anda baik-baik saja.

\begin{enumerate}
  \item Jika hitungan mundur berakhir dan Anda tidak merespons, secara otomatis memanggil telepon darurat.
  \item Mengirim lokasi Anda ke kontak darurat.
  \item Model yang Didukung:
\end{enumerate}

\booksubsection{Apple}

\begin{itemize}
  \item \textbf{Apple}Android
  \item \textbf{Android}Catatan:
\end{itemize}

\booksubsection{Fitur harus diaktifkan terlebih dahulu untuk dapat digunakan.}

\begin{itemize}
  \item Dapat terpicu secara salah (misalnya roller coaster, olahraga intens), cukup batalkan saja.
  \item Bergantung pada sensor dan algoritma, tidak 100\% andal, namun menambah satu lapis perlindungan ekstra.
  \item Empat, Teknologi SparkLink: Koneksi Nirkabel Generasi Berikutnya
\end{itemize}

\booksection{SparkLink adalah teknologi komunikasi nirkabel generasi baru yang dipimpin pengembangannya oleh Huawei, dengan target menggantikan Bluetooth dan Wi-Fi dalam beberapa skenario. Pada indikator kunci, ia menandingi atau bahkan melampaui solusi yang ada: kecepatan air interface dapat mencapai 900Mbps (jauh melampaui 50Mbps Bluetooth), latensi serendah 20 $\mu$ s (Bluetooth sekitar 15ms), satu jaringan mendukung 4096 perangkat terhubung secara bersamaan (Bluetooth biasanya di bawah 8).}

Karakteristik Teknologi:

\booksubsection{Latensi ultra rendah}

\begin{itemize}
  \item \textbf{: Latensi end-to-end hanya 20 mikrodetik (0.02ms), dua tingkat lebih rendah dari Bluetooth, sinkronisasi game dan musik hampir tidak terasa.}Kecepatan tinggi
  \item \textbf{: Kecepatan puncak single-link 900Mbps, puluhan kali lipat audio Bluetooth, dapat mentransmisikan audio lossless bahkan video HD.}Konsumsi daya rendah
  \item \textbf{: Konsumsi daya standby setara dengan Bluetooth, namun efisiensi transmisi lebih tinggi, konsumsi energi per unit data lebih rendah.}Cakupan luas + konkurensi tinggi
  \item \textbf{: Jarak koneksi 50\% lebih jauh dari Bluetooth, satu jaringan mendukung ribuan perangkat online secara bersamaan.}Skenario Aplikasi:
\end{itemize}

\booksubsection{Headphone nirkabel}

\begin{itemize}
  \item \textbf{: Latensi lebih rendah, saat bermain game sinkronisasi audio dan visual.}Jam tangan pintar
  \item \textbf{: Koneksi lebih stabil.}Rumah pintar
  \item \textbf{Rumah pintar}Interkoneksi otomotif
  \item \textbf{: Koneksi ponsel-head unit yang lebih mulus.}Perangkat yang Didukung:
\end{itemize}

\booksubsection{Saat ini terutama perangkat ekosistem Huawei, seperti seri Huawei Mate 60, sebagian headphone dan jam tangan Huawei. Di masa depan akan ada lebih banyak merek yang mendukung.}

Prospek Teknologi

\textbf{Teknologi SparkLink mewakili arah masa depan komunikasi nirkabel. Seiring bergabungnya lebih banyak produsen, kita mungkin akan melihat lebih banyak perangkat yang mendukung SparkLink, menikmati pengalaman nirkabel yang lebih baik.}\par

Periksa: Apakah ponsel Anda mendukung fitur-fitur ini

\begin{tcbpractice}{Cari "satelit" di pengaturan, lihat apakah ada opsi terkait komunikasi satelit.}
\begin{enumerate}
  \item Cari "deteksi kecelakaan" atau "SOS darurat", periksa apakah diaktifkan.
  \item Jika ponsel Anda adalah Huawei, lihat apakah mendukung koneksi perangkat SparkLink.
  \item Pelajari cara penggunaan fitur-fitur ini, mungkin berguna di saat kritis.
  \item II. Bagian Perangkat Lunak (Bab 8-14)
\end{enumerate}

\end{tcbpractice}
